\section{北线不是一条箭头：蜀汉怎样反复抵达陇右}

\subsection{汉中先是一项守备，再是一项进攻}

诸葛亮驻汉中时，成都面对的不是一张可以随意向北延长的地图。汉中盆地、褒斜
与祁山方向的道路，把粮食、马匹、传令者和守军分在一连串并不相等的节点上。
《三国志·诸葛亮传》把北伐的诏令、进退和将领置于清晰的年序中；这正是理解
行动的骨架。可是年序并不能取消道路：同一份动员一离开成都，便要由仓口、山
谷和地方人力把它分解成许多可失败的小事。

《出师表》把益州的疲弊写进出兵理由。它不是后方账本，也不是诸葛亮的内心
独白；它是一份向君主陈述秩序、任用与北向选择的政治文字。读它时应当看见
成都为何要把有限资源压向汉中，同时也应避免从“疲弊”推出精确户数、粮额或
每一位官员的态度。文字使动员具有可公开辩护的语言，却不能替运输者跨过山口。

二二八年祁山方向的响应一度改变了陇右的局面。南安、天水、安定的消息对蜀军
很重要，因为战争并非只在两支主力之间发生；城郡是否转向，决定了军队能否
取得粮食、向导、情报和暂时的立足处。街亭失守使这一段联系收缩。将马谡的
处置当作故事终点，会遮蔽更难的事实：蜀军缺少能够替换的安全补给线，一处
关键据点不能维持，就会使前方的响应成为难以守住的负担。撤军是承认这种关系
已不能继续承受，而非一句道德判断可以解释的结局。

\subsection{祁山、武都与阴平：小国的边缘不是固定边界}

二二九年取得武都、阴平，在蜀汉叙事中往往像街亭之后的一次修补。若从汉中看，
它的意义更具体：向西北推出一层可供警戒和周旋的空间，或可降低敌军直接逼近
山口的速度。它没有把蜀军变成陇右的永久主人。新的驻点需要补给，需要地方
消息，也会把守备人数从其他方向抽走。所谓“得郡”是一项对距离的暂时处理，
不是一笔已经兑现的领土收入。

《华阳国志》保存蜀地的地方传统和后出的叙事线索，可以帮助读者不把成都以外
写成空白；但它的成书时间和立场也要求我们谨慎。它不能单独为每一处山道和
每一个地方首领提供即时证词。把它与《三国志》及历史地理研究并读，能较稳妥
地看见盆地、汉中和陇右之间的不同空间，却不能为蜀汉画出精确而静止的北界。

二三一年李严因转运问题被废，正好暴露出国家想把哪一种风险归入责任。史书
留给我们的有诏令、弹劾和处分，也有后来围绕人物品行的叙述。无论谁应承担
多少责任，山道上的迟滞都不可能只由一个官员创造或消除。成都需要有人保证
前线供应，前线需要能理解季节、道路和仓储的执行者；当承诺落空，朝廷必须
以处分重新宣布命令仍有边界。制度由此得到可追责的入口，代价则是复杂的运输
困境被压缩为一个人是否称职的判断。

\subsection{五丈原的时间}

五丈原并不要求读者相信一场决定天下的正面对决。诸葛亮在这里选择屯驻，说明
蜀军试图让前线生产和驻守减轻远程运输的压力；司马懿的相持，则让时间成为
另一种武器。两种选择都被资源限定：蜀军无法无限停留，魏军也要维持关中防线
和其他方向的调度。诸葛亮病逝、军队撤回是可确认的顺序，营垒大小、兵粮数字
和每一次军议却不宜以传说填补。

撤军之后，蒋琬接任并不只是“守成”。他必须处理军队回到汉中后的安置、成都
官署的连续性、与吴的关系以及北线仍需守住的山口。费祎后来限制大举北伐，也
是在同一组选择中分配人力。若把谨慎只写作缺乏胆略，便会把蜀汉唯一能够长久
存在的条件抹去：一个人口与道路都较受限制的政权，必须允许一些资源留在没有
战功名称的守备、仓储和地方协调之中。

\subsection{姜维的机会与反作用}

姜维北向时，陇西羌胡的变动有时提供机会。边地集团并非蜀魏两国任意推动的
棋子；他们有自己的迁徙、联盟和生计，现存汉文史料大多从蜀魏军政的视角记录
接触。可以确认的出兵与撤退，不能替他们写出统一的政治意图。正因如此，姜维
试图利用边地局势时，既可能迫使魏军分心，也必须依赖无法完全控制的地方条件。

二六二年狄道之役的撤退，把这个反作用再次显露出来。邓艾据险来援、蜀军粮运
受压，是材料能够支持的行动骨架；没有连续账簿，不能把双方消耗换算成精确
数字。它也不单独“造成”二六三年的灭蜀。更合适的理解是：每一次无法把机会
转成可持续驻守的北向行动，都会让成都的可用余地更窄，使下一次面对魏军多路
南下时，汉中与剑阁之外能够临时调动的资源更少。
